\documentclass[10pt,twocolumn]{article}
\usepackage[top=0.97in,bottom=0.97in,lmargin=0.75in,rmargin=0.75in]{geometry}
\usepackage[labelfont=bf]{caption}

\setlength{\columnsep}{2.5em}


\setlength{\columnsep}{2.5em}

\usepackage{mathpazo,helvet,inconsolata,bbding}

\usepackage[utf8]{inputenc}
\usepackage[autolanguage,np]{numprint}

%\usepackage{nopageno}
\usepackage[english]{babel}
\usepackage{csquotes}
\usepackage{graphicx}
\usepackage{textcomp,gensymb}
\usepackage{amsmath}
\usepackage{amssymb}
%\usepackage{algorithm}
%\usepackage{algpseudocode}
\usepackage{xspace}
%\usepackage{relsize}
\usepackage{stackengine}
\usepackage{tabularx}
%\usepackage{pgfplots}
\usepackage{authblk}
%\pgfplotsset{width=10cm,compat=1.9}
%\usepgfplotslibrary{external}
%\tikzexternalize

\usepackage[
style=authoryear-comp,
giveninits=true,
uniquename=init,
citestyle=authoryear,
natbib=true,
%backend=biber
]{biblatex}
\bibliography{ref,ref2}

% \usepackage[
% style=numeric,
% sorting=none,
% style=authoryear-icomp,
% giveninits=true,
% uniquename=init,
% citestyle=numeric,
% natbib=true,
% uniquelist=false,
% isbn=false
% ]{biblatex}

\bibliography{ref}
\renewcommand*{\bibfont}{\scriptsize}
\AtEveryBibitem{\clearlist{language}}
\AtEveryBibitem{\clearfield{url}\clearfield{urlyear}}
\AtEveryBibitem{\iffieldundef{doi}{}{\clearfield{url}\clearfield{urlyear}}}
\AtEveryBibitem{\clearfield{month}}
\AtEveryBibitem{\clearlist{publisher}}
\AtEveryBibitem{\clearfield{note}}
\AtEveryBibitem{\clearname{editor}}
\AtEveryBibitem{\clearfield{series}}
\AtEveryBibitem{\clearfield{issn}}
\AtEveryCitekey{\clearfield{issn}}
\renewbibmacro*{url+urldate}{}

\stackMath
\setstackgap{S}{2pt}

\DeclareMathOperator*{\argmax}{arg\,max}
\DeclareMathOperator*{\argmin}{arg\,min}
\DeclareMathOperator*{\argminA}{arg\,min} % Jan Hlavacek
\DeclareMathOperator*{\fr}{\overrightarrow{r}}
\DeclareMathOperator*{\frp}{{\overrightarrow{r}}'}
\DeclareMathOperator*{\fmu}{\overrightarrow{\mu}}
\DeclareMathOperator*{\E}{\mathbb{E}}
\DeclareMathOperator*{\Agg}{\sum_{j=1}^{N}}
\DeclareMathOperator*{\Gagg}{\bigodot}
\DeclareMathOperator*{\gnn}{GNN}
\DeclareMathOperator*{\inr}{INR}
\DeclareMathOperator*{\mlp}{MLP}
\DeclareMathOperator*{\relu}{ReLU}
\newcommand{\dt}{\Delta t \:}

\newcommand{\force}{\textrm{Force}}
\newcommand{\loss}{L}

\renewcommand\Vec[1]{\ensuremath{\boldsymbol{#1}}}
\newcommand\Mat[1]{\ensuremath{\boldsymbol{#1}}}
\newcommand\trnsp{\ensuremath{\mathsf{T}}}

\usepackage{booktabs}
\usepackage{pifont}
\usepackage{xcolor}
\newcommand{\cmark}{\textcolor{green!60!black}{\ding{51}}}
\newcommand{\xmark}{\textcolor{red!70!black}{\ding{55}}}
\newcommand{\wmark}{\textcolor{orange!80!black}{$\circ$}}

\usepackage[dvipsnames]{xcolor}
\usepackage{hyperref}
\hypersetup{%
	colorlinks=true,%
	linkcolor=BrickRed,%
	citecolor=ForestGreen,
	urlcolor=NavyBlue
}
\usepackage[capitalize,noabbrev,nameinlink]{cleveref}
\crefname{figure}{Figure}{Figures}
\crefname{appendixfigure}{Supplementary Figure}{Supplementary Figures}
\crefname{table}{Table}{Tables}
\crefname{appendixtable}{Supplementary Table}{Supplementary Tables}
\crefname{equation}{Equation}{Equations}
\crefname{appendixequation}{Supplementary Equation}{Supplementary Equations}
\creflabelformat{equation}{#2#1#3}
\creflabelformat{appendixequation}{#2#1#3}

\usepackage{caption}
\captionsetup{font=footnotesize}

\def\etc{etc.\@\xspace}
\def\eg{e.g.\@\xspace}
\def\Etc{Etc.\@\xspace}
\def\Eg{E.g.\@\xspace}

\setlength{\fboxsep}{0pt} % Adjust the padding
\setlength{\fboxrule}{0.25pt} % Adjust the border width
\setlength{\belowcaptionskip}{-2pt}



\newcommand{\ftile}[1]{\raisebox{-\height+2ex}{\fbox{\includegraphics[width=\linewidth]{#1}}}}
%\title{Graph neural networks can uncover the hidden structure and function underlying the observed temporal activity in heterogeneous neural assemblies}
\title{Experiment-LLM-memory}
\author{C\'edric Allier}
\author{Stephan Saalfeld $^\text{\Envelope}$}
\affil{Janelia Research Campus, Howard Hughes Medical Institute, Ashburn, VA 20147, USA.}
\affil{\href{mailto:saalfelds@janelia.hhmi.org}{\small\Envelope~saalfelds@janelia.hhmi.org}}
\date{\today}


\begin{document}

\twocolumn[ 
\begin{@twocolumnfalse}
\maketitle


\begin{abstract}
\noindent{} We previously showed that graph neural networks can recover circuit structure and signaling functions of neural assemblies from activity data alone. However, this inverse problem is known to be ill-posed under certain conditions. For instance, when neural activity is low-rank, many different circuits can generate the same neuron traces. It remains an open question whether graph neural networks can recover connectivity in general, or only for specific classes of neural assemblies. To address this question, we extend our graph neural network framework with a large language model. In this new setup, the role of the graph neural network framework is to perform experiments and deliver quantified results. The role of the large language model is to interpret, compare and finally compress the experiment results into an structured memory. On the basis thereof, the large language model is next prompted to mutate selectively the neural dynamics configuration and/or the graph neural network training scheme. We show that these sequential interactions between experimentation, large language model and long-term memory, leads progressively to a scientific tool. Testable hypotheses are drawn, repeatable experiments are conducted to validate or falsify them, and ultimately causal understanding emerge. Importantly, the closed-loop scientific reasoning results from the interactions between the three components rather than residing solely within the large language model.

\end{abstract}


\vspace{3em}

\end{@twocolumnfalse} 
]

\section{Introduction}

Understanding when graph neural networks (GNNs) can recover synaptic connectivity from neural activity, and when they fundamentally fail, requires exploring a high-dimensional landscape of neural activity regimes. A classical approach would rely on nested grid searches: first over classes of neural assemblies to generate activity, and then over GNN training hyperparameters for each regime. Such an approach rapidly becomes intractable, prone to gaps, and poorly suited to uncovering general principles across regimes.

\paragraph{LLM-guided discovery.}
Recent methods coupling large language models (LLMs) with external validation, including FunSearch~\citep{romera-paredes_mathematical_2024}, AlphaEvolve~\citep{novikov_alphaevolve_2025}, OPRO~\citep{yang_large_2024}, and Google Tree Search~\citep{aygun_ai_2025}, have demonstrated the ability to explore problem-solution landscape in mathematics and computer science. These systems iteratively select candidate solutions based on solution--score tuples and use LLMs to propose refinements or fill gaps in the search space. While highly effective for optimization, these approaches are however limited by the finite context window of LLMs. Without explicit long-term memory, hypotheses and deductions can be discarded, preventing accumulation of understanding.
Other systems combining LLMs with external experimentation or validation are the AI Scientist~\citep{lu_ai_2024} and the Virtual Lab~\citep{swanson_virtual_2024}. The former automates machine learning research; the latter coordinates wet-lab experiments. However, both treat experiments as largely independent, lacking mechanisms to accumulate mechanistic understanding across runs. Knowledge does not persist or transfer beyond local optimization.

Conversely, approaches such as Titans~\citep{behrouz_titans_2024} and Anthropic's engineering guidelines~\citep{anthropic2025harnesses} introduce external memory to overcome the LLM's context window limitation. Titans implements memory within neural networks trained by gradient descent, while Anthropic advocates text-based progress files to ensure consistency across LLM outputs. In both cases, memory supports coherence, but without coupling to external experiments, generated insights cannot be empirically tested or refuted, leaving their validity uncertain. 
Alternatively, other models such as OpenAI’s o1/o3~\citep{openai2024reasoning, openai2025o3mini} and DeepSeek’s R1~\citep{deepseek-ai_deepseek-r1_2025} extend LLM reasoning by scaling test-time computation and structuring deliberation into long chains of thought, often coupled to code execution or proof verification. These models achieve state-of-the-art performance on many benchmarks~\citep{li_system_2025}. But knowledge accumulation is only possible with an external observer. Without long-term memory, it can not be internalized by the system itself.

Other systems such as Robin~\citep{ghareeb_robin_2025}, AI co-scientist~\citep{gottweis_towards_2025}, and Kosmos~\citep{mitchener_kosmos_2025} couple LLMs to large-scale literature analysis and world models to generate hypotheses. This approach can perform extensive synthesis across existing datasets and publications, enabling deep deductive insights. However, when hypotheses are derived from literature alone, the process can be short-ended when falsification needs additional studies.

Finally, recent works on embodied agents do integrate experimentation, LLM and structured memory, enabling trial-and-error learning but not hypothesis-driven falsification. As such, these works target task completion rather than scientific understanding~\citep{majumder_clin_2023, wang_voyager_2023, yang_agentic_2025}.


\paragraph{The gap.}
Here, we show that accumulating scientific understanding can be accomplished by the coupling of three components: external experiments, LLM-based hypothesis generation, and explicit long-term memory. We present a system that closes this loop. A graph neural network framework performs controlled experiments on simulated neural dynamics and produces quantitative evaluations. A large language model interprets these results, compresses them into structured long-term memory, and intervenes by mutating either the neural dynamics configuration or the GNN training scheme. These interventions are associated with predictions, which are tested by subsequent experiments. Memory is updated accordingly: reinforced when predictions are confirmed and revised when falsified.

We demonstrate this framework on the problem of neural connectivity inference: determining when GNNs can recover the synaptic weight matrix $W_{ij}$ from neural activity $\{x_i(t)\}$. Across 2048 GNN optimizations spanning 128 neural activity regimes, the system accumulates explicit transferable understanding. Induction identifies regularities across regimes, such as scaling relationships between learning rate ratios and constraint complexity. Abduction proposes explanatory mechanisms, for example that data effective rank determines the attainable connectivity $R^2$. Deduction generates testable predictions, such as increasing $n_{\mathrm{frames}}$ to raise effective rank and enable recovery. Falsification rejects or refines these predictions when outcomes contradict expectations. Through these processes, experimental results are consolidated into explicit memory as hypotheses and regime-level knowledge that persist across blocks and systematically constrain, inform, and accelerate future exploration.


\section{Methods}

\paragraph{Experimental setup.}
We employ the GNN framework introduced in [...] to model neural dynamics with graph-computation (appendix~\ref{app:neural_dynamics}): neuron-neuron interactions are computed first, then aggregated to update neuron states. The computation implemented with the PyTorch Geometric library \citep{fey_fast_2019} is differentiable, so interaction functions, neuron properties, and update rules can be learned from data. This poses an inverse problem over the space $\mathcal{G}$ of graph-computations modeling neural dynamics: given activity traces $\{x_i(t)\}$ produced by some $g \in \mathcal{G}$, recover $g$. In this experimental setup we address the existence and uniqueness of solutions obtained through GNN optimization.

\paragraph{Experiments - LLM - memory.}
The system couples the three components, i.e. external experiments, LLM reasoning, and long-term memory, through simple text files exchange. Two text files structure the LLM's scientific process. The \emph{instruction file} (appendix~\ref{app:instruction}) defines the inverse problem to solve and the methodology to follow: parameter ranges to explore through tree search, convergence criteria, and logging format. It also contains a brief domain knowledge: GNN structure and linear algebra principles. The \emph{memory file} (appendix~\ref{app:memory}) maintains the LLM's accumulated understanding in two tiers: a set of established principles and open questions that grows across experimentation, and a working memory of current hypotheses and recent iterations. This two-tier structure enables both long-term knowledge accumulation and short-term memory. The LLM interact with the experiments through two other text files. It can modify the configuration file specifying simulation parameters (connectivity type, noise, $n_{\mathrm{neurons}}$) and GNN parameters (learning rates, architecture, regularization). When the experiment is executed it writes the GNN results into a text log-file which, serves as an episodic memory for the LLM.  
\\
Experiments are organized into \emph{blocks}. Between each block, the LLM fixes a neural assembly $g \in \mathcal{G}$. Within a block, the LLM explores GNN training parameters (learning rates, regularization) over 16 iterations. At each iteration, a different activity traces $\{x_i(t)\}$ is generated from $g$, and the GNN is optimized with new parameters set by the LLM. Parameter exploration follows Upper Confidence Bound (UCB) tree search, adapted from \citet{romera-paredes_mathematical_2024}. Each configuration is a node; the LLM selects the next parent by balancing exploitation (high-performing) with exploration (under-visited branches) via $\mathrm{UCB}(k) = R^2_k + c\bigl(\ln N / n_k\bigr)^{1/2}$, where $R^2_k$ is the mean coefficient of determination between true and learned connectivity weights $W_{ij}$.This metric scores model recovery at node $k$ of search tree, $n_k$ its visit count, and $N$ the total iterations. The LLM mutates GNN training parameters from the selected parent and logs its reasoning.

At block boundaries, the LLM performs three operations. First, to overcome rigid rule-based exploration, it evaluates its tree search strategy (branching rate, improvement rate, stuck detection) and edits the \emph{instruction} file itself to add or modify rules. Second, it selects the next simulation regime to address gaps in the search space or refine optimization of a previous block. Third, it compresses the block into the knowledge base: a neural activity regime comparison table (best $R^2$, optimal parameters, key findings), established principles that transfer across neural activity regimes, and open questions for future blocks.


\paragraph{Epistemic metrics and ablation studies.}
We introduce \emph{epistemic metrics} (Appendix~\ref{appendix:epistemic_metrics}) to e.g. quantify how the system acquires, tests, revises, and transfers knowledge across experimental blocks. Importantly, these metrics do not measure task performance, but the structure and dynamics of the learning process. 

To isolate the contribution of each component to these epistemic behaviors, we perform controlled ablation studies in which the external experimental engine and evaluation pipeline are held fixed, while selectively modifying the decision module. Specifically, we consider: (i) a \emph{no-memory} ablation, in which the large language model receives only the last $K$ iterations and has no access to accumulated long-term knowledge; (ii) a \emph{no-UCB-tree} ablation, in which parameter mutations are sampled uniformly or from matched heuristic distributions rather than guided by UCB exploration; and (iii) a \emph{frozen-protocol} ablation, in which the large language model may propose parameter updates but is prohibited from editing the exploration rules. These ablations allow us to attribute observed epistemic behaviors to the interaction between persistent memory, guided exploration, and adaptive protocol control.





% Analysis.log is episodic memory
% Self-mutation to overcome rigid rule-based system, dynamics demand continuous learning and adaptation as RL
% No need to give extended logical instructions, the LLM apparently already encodes all necessary rules.
% Monte-Carlo tree search, UCB-tree to guide exploration.
% Literature injection
% Code mutation to extend the framework from parameter to parameter/code exploration  








\section{Results}

\begin{figure*}[htbp]
\centering
\includegraphics[width=\textwidth]{signal_chaotic_1_Claude_epistemic_timeline.png}
\caption{\textbf{Epistemic reasoning timeline across 149 iterations (11 blocks).} Each node represents a reasoning event, colored by mode and sized by significance. Blocks alternate white/gray backgrounds. Edges show causal relationships: solid gray (leads to), dashed blue (triggers), red backward arrows (to Falsification). \textbf{Modes:} \emph{Induction}: pattern extraction from observations; \emph{Boundary}: parameter limit testing; \emph{Abduction}: causal hypothesis generation; \emph{Deduction}: prediction from hypothesis; \emph{Falsification}: hypothesis rejection; \emph{Analogy/Transfer}: cross-regime knowledge transfer; \emph{Meta-reasoning}: strategy adaptation; \emph{Regime}: phase identification; \emph{Uncertainty}: stochasticity awareness; \emph{Constraint}: parameter constraint formulation; \emph{Predictive}: quantitative rule derivation; \emph{Causal}: mechanistic chain modeling.}

\label{fig:epistemic_timeline}
\end{figure*}

We found a pyramidal reasoning scheme as shown in \cref{fig:epistemic_timeline}. At the bottom, inductions and abductions occur throughout the process, forming the foundation of hypothesis generation and pattern recognition. This layer supports a second layer of deductions, testable predictions derived from hypotheses. All of these reasoning events can be falsified or refined, as shown by the red nodes appearing after blue predictions. Ultimately, causal chains of reasoning emerge. The first appeared at iteration 48, where the chain $\texttt{connectivity\_rank} \to \texttt{eff\_rank} \to R^2_{\text{ceiling}} \to \texttt{convergence}$ was established when the LLM constructed a multi-step mechanistic model explaining why low-rank connectivity limits GNN recovery. A second causal node at iteration 127 formalized the $\texttt{low\_rank=10} \to \texttt{eff\_rank=4} \to \texttt{unlearnable}$ barrier. At iteration 135, Dale's law stabilization ($\rho: 1.0 \to 0.65$) was causally linked to increased effective rank. Cross-block edges (dashed blue spanning blocks) show knowledge transfer, with block summaries triggering analogical reasoning in subsequent regimes.

Key findings include: a Hypothesis Turnover Rate of 2.43 indicating iterative refinement cycles; Deductive Accuracy of 71\% reflecting calibrated prediction; Transfer Success Rate of only 35\% highlighting the challenge of cross-regime generalization; and a Causal Depth of 3.6 demonstrating multi-step mechanistic reasoning. The first causal chain emerged at iteration 48, marking the transition from parameter search to explanatory modeling.


\section{Discussion}

\citep{collins_building_2024} critique pure LLM scaling and advocate for structured world models. They focus on human-LLM collaboration through mutual inference, but the LLM assists while the human drives the scientific process. We show how to elevate the LLM to equal partnership, where both human and LLM hypothesize, intervene, and falsify.

Symbolic logic seems to be structurally encoded in LLM. It appears that the LLM coupled to experimentation and long-term memory exhibits rapidly a tropism for logical reasoning. The current approach has thus a key strength: the epistemic reasoning emerges naturally from the interaction between the system components, it's not prescribed or forced. We did not injected explicit epistemic tags into the instruction file.  

We summarize the distinction between LLM-guided approaches upon a new perspective which could serve as guidelines for AI-scientific tool. In assessing LLM-guided scientific systems, we consider the availability of three components (external experiments, LLM, explicit long-term memory) and the capacity for four types of reasoning (induction, abduction, deduction, falsification).

The current system demonstrates the methodology on a single domain (neural dynamics). Yet the approach can be extended beyond this field. Any domain expressible as local interactions aggregated over a graph, e.g. molecular systems, traffic networks, social dynamics, gene regulatory circuits, fluid discretizations, falls within the proposed framework. Scaling from many regimes to many domains is therefore a natural extension. Work on simulation is a necessary prerequisite to real data. For each simulation domain, establishing learnability and intelligibility boundaries is a generic problem. Under what conditions can parameters be recovered (learnability), and under what conditions can the recovery be attributed to mechanistic understanding rather than statistical fitting (intelligibility) ? 

% What will not be convincing to reviewers
% “The LLM reasons” (too easy to attack philosophically).
% Improvements in final score alone (could just be UCB + luck + optimization).
% Long chains of explanation in text (people will say it’s post-hoc narration).
% The risk is reviewers saying: “This is just hyperparameter optimization with an LLM in the loop.”
% You beat that by making the paper about epistemic behavior, not optimization:
% diagnostic experiments
% prediction + falsification
% explicit strategy edits that improve sample-efficiency
% regime classification in memory that generalizes


\section{Conclusion}


\section*{Acknowledgements}


\printbibliography

\appendix
\onecolumn
\clearpage

\onecolumn

\appendix



\section{Neural dynamics and GNN architecture}
\label{app:neural_dynamics}

\subsection*{Neural dynamics}

Neural assemblies are simulated as graphs of $N$ neurons with weighted synaptic connections $W_{ij}$. Activity $x_i$ evolves according to:
\begin{equation}
\dot{x}_i = -\frac{x_i}{\tau_i} + s_i\phi(x_i) + g_i\sum_j W_{ij} \psi(x_j) + \eta_i(t)
\end{equation}
where $\tau_i$ controls decay, $s_i$ self-coupling, $g_i$ scales aggregated inputs, and $\eta_i(t)$ is Gaussian noise. Weights are drawn from .... Transfer functions $\phi, \psi$ are $\tanh$ variants; neuron types are parameterized by ... $.

\subsection*{GNN Architecture}

The GNN learns:
\begin{equation}
\widehat{\dot{x}}_i = \phi^*(\mathbf{a}_i, x_i) + \sum_j W_{ij} \psi^*(\mathbf{a}_j, x_j)
\end{equation}
where $\phi^*, \psi^*$ are MLPs (ReLU, 3 layers, hidden dim 64), $\mathbf{a}_i \in \mathbb{R}^2$ are learnable neuron embeddings, and $W_{ij}$ is the learnable connectivity matrix. The loss combines prediction error with regularization terms enforcing zero steady state, exponential decay, monotonic transfer functions, and weight sparsity.

\section{Instruction File}
\label{app:instruction}

\begin{small}
\begin{verbatim}
# Simulation-GNN Training Landscape Study

## Goal

Map the **simulation-GNN training landscape**: understand which neural activity simulation configurations allow successful GNN training (connectivity_R2 > 0.9) and which are fundamentally harder.

## Iteration Loop Structure

Each block = `n_iter_block` iterations exploring one simulation configuration.
The prompt provides: `Block info: block {block_number}, iteration {iter_in_block}/{n_iter_block} within block`

## File Structure (CRITICAL)

You maintain **TWO** files:

### 1. Full Log (append-only record)

**File**: `{config}_analysis.md`

- Append every iteration's full log entry
- Append block summaries
- **Never read this file** — it's for human record only

### 2. Working Memory

**File**: `{config}_memory.md`

- **READ at start of each iteration**
- **UPDATE at end of each iteration**
- Contains: established principles + previous blocks summary + current block iterations
- Fixed size (~500 lines max)

---

## Iteration Workflow (Steps 1-5, every iteration)

### Step 1: Read Working Memory

Read `{config}_memory.md` to recall:

- Established principles
- Previous block findings
- Current block progress

### Step 2: Analyze Current Results

**Metrics from `analysis.log`:**

- `spectral_radius`: eigenvalue analysis of connectivity
- `svd_rank`: SVD rank at 99% variance (activity complexity)
- `test_R2`: R² between ground truth and rollout prediction
- `test_pearson`: Pearson correlation between ground truth and rollout prediction
- `connectivity_R2`: R² of learned vs true connectivity weights
- `cluster_accuracy`: clustering accuracy (neuron type classification)
- `final_loss`: final training loss

**Classification:**

- **Converged**: connectivity_R2 > 0.9
- **Partial**: connectivity_R2 0.1-0.9
- **Failed**: connectivity_R2 < 0.1

**UCB scores from `ucb_scores.txt`:**

- Provides computed UCB scores for all exploration nodes including current iteration
- At block boundaries, the UCB file will be empty (erased). When empty, use `parent=root`

Example:

```
Node 2: UCB=2.175, parent=1, visits=1, R2=0.997
Node 1: UCB=2.110, parent=root, visits=2, R2=0.934
```

### Step 3: Write Outputs

Append to Full Log (`{config}_analysis.md`) and **Current Block** sections of `{config}_memory.md`:

- In memory.md: Insert iteration log in "Iterations This Block" section (BEFORE "Emerging Observations")
- Update "Emerging Observations" at the END of the file

**Log Form:**

```
## Iter N: [converged/partial/failed]
Node: id=N, parent=P
Mode/Strategy: [success-exploit/failure-probe]/[exploit/explore/boundary]
Config: lr_W=X, lr=Y, lr_emb=Z, coeff_W_L1=W, batch_size=B, low_rank_factorization=[T/F], low_rank=R, n_frames=NF
Metrics: test_R2=A, test_pearson=B, connectivity_R2=C, cluster_accuracy=D, final_loss=E
Activity: [brief description]
Mutation: [param]: [old] -> [new]
Parent rule: [one line]
Observation: [one line]
Next: parent=P
```

### Step 4: Parent Selection Rule in UCB tree

Step A: Select parent node

- Read `ucb_scores.txt`
- If empty → `parent=root`
- Otherwise → select node with **highest UCB** as parent

**CRITICAL**: The `parent=P` in the Node line must be the **node ID** (integer) of the selected parent, NOT "root" (unless UCB file is empty). Example: if you select node 3 as parent, write `Node: id=4, parent=3`.

Step B: Choose strategy

| Condition                            | Strategy             | Action                                       |
| ------------------------------------ | -------------------- | -------------------------------------------- |
| Default                              | **exploit**          | Highest UCB node, try mutation               |
| 3+ consecutive R² ≥ 0.9              | **failure-probe**    | Extreme parameter to find boundary           |
| n_iter_block/4 consecutive successes | **explore**          | Select outside recent chain                  |
| Good config found                    | **robustness-test**  | Re-run same config                           |
| 2+ distant nodes with R² > 0.9       | **recombine**        | Merge params from both nodes                 |
| 100% convergence, branching<10%      | **forced-branch**    | Select node in bottom 50% of tree            |
| 4+ consecutive same-param mutations  | **switch-dimension** | Mutate different parameter than recent chain |
| 3+ partial results probing boundary  | **boundary-skip**    | Accept boundary as found, explore elsewhere  |
| 8+ consecutive sequential (no branch)| **forced-diversity** | Select any node with visits ≥ 3 that is NOT the most recent 4 nodes |

**Recombination details:**

Trigger: exists Node A and Node B where:

- Both R² > 0.9
- Not parent-child (distance ≥ 2 in tree)
- Different parameter strengths

Action:

- parent = higher R² node
- Mutation = adopt best param from other node

Example:

```
Node 12: lr_W=1E-2, lr=1E-4, R²=0.94  (good lr_W)
Node 38: lr_W=5E-3, lr=2E-3, R²=0.97  (good lr)

Recombine → lr_W=1E-2, lr=2E-3
```

### Step 5: Edit Config File

Edit config file for next iteration of the exploration.
(The config path is provided in the prompt as "Current config")

**CRITICAL: Config Parameter Constraints**

**DO NOT add new parameters to the `claude:` section.** Only these fields are allowed:

- `n_epochs`: int (training epochs per iteration)
- `data_augmentation_loop`: int (data augmentation count)
- `n_iter_block`: int (iterations per block)
- `ucb_c`: float value (0.5-3.0)

Any other parameters belong in the `training:` or `simulation:` sections, NOT in `claude:`.

Adding invalid parameters to `claude:` will cause a validation error and crash the experiment.

**Training Parameters (change within block):**

```yaml
training:
  learning_rate_W_start: 2.0E-3 # range: 1E-4 to 1E-2
  learning_rate_start: 1.0E-4 # range: 1E-5 to 1E-3
  learning_rate_embedding_start: 2.5E-4 # used when n_neuron_types>1
  coeff_W_L1: 1.0E-5 # range: 1E-6 to 1E-3
  batch_size: 8 # values: 8, 16, 32
  low_rank_factorization: False # or True
  low_rank: 20 # range: 5-100
  coeff_edge_diff: 100 # enforces positive monotonicity
```

**Simulation Parameters (change at block boundaries only):**

```yaml
simulation:
  n_frames: 10000
  connectivity_type: "chaotic" # or "low_rank"
  Dale_law: False # or True
  Dale_law_factor: 0.5
  connectivity_rank: 20 # if low_rank
  n_neurons: 100 # can be changed to 1000 and ONLY 1000, if Iter > 512
  n_neuron_types: 1 #  # can be changed to 4 and ONLY 4,  if Iter > 1024
```

**Claude Exploration Parameters:**

These parameters control the UCB exploration strategy. Can be adjusted between blocks to adapt exploration behavior.

```yaml
claude:
  ucb_c: 1.414 # UCB exploration constant (0.5-3.0), adjust between blocks
```

**UCB exploration constant (ucb_c):**

- `ucb_c` controls exploration vs exploitation: UCB(k) = R²_k + c × sqrt(ln(N) / n_k)
- Higher c (>1.5) → more exploration of under-visited branches
- Lower c (<1.0) → more exploitation of high-performing nodes
- Default: 1.414 (√2, standard UCB1)
- Adjust between blocks based on search behavior:
  - If stuck in local optimum (all R² similar, no improvement) → INCREASE ucb_c to 2.0
  - If too much random exploration (jumping between distant nodes) → DECREASE ucb_c to 1.0
  - Typical range: 0.5 to 3.0

---

## Block Workflow (Steps 1-3, every end of block)

Triggered when `iter_in_block == n_iter_block`

### STEP 1: COMPULSORY — Edit Instructions (this file)

You \***\*MUST\*\*** use the Edit tool to add/modify parent selection rules in this file.
Do NOT just write recommendations in the analysis log — actually edit the file.

After editing, state in analysis log: `"INSTRUCTIONS EDITED: added rule [X]"` or `"INSTRUCTIONS EDITED: modified [Y]"`

**Evaluate and modify rules based on:**

**Branching rate:**

- Branching rate < 20% → ADD exploration rule
- Branching rate 20-80% → No change needed
- **Branch rate** = parent ≠ (current_iter - 1), calculate for **entire block**

Example:

```
Block 1 (16 iterations):
Sequential: Iters 1-14 (all parent = node-1)
Branches: Iter 15 (parent=2)
Branches: 1 out of 15 → Branching rate = 7%
```

**Improvement rate:**

- If <30% improving → INCREASE exploitation (raise R² threshold)
- If >80% improving → INCREASE exploration (probe boundaries)

**Stuck detection:**

- Same R² plateau (±0.05) for 3+ iters? → ADD forced branching rule

**Dimension diversity:**

- Count consecutive iterations mutating **same parameter**
- If > 4 consecutive same-param → ADD switch-dimension rule

Example:

```
Iter 2-14: all mutated lr_W → 13 consecutive same-dimension → ADD rule
```

### STEP 2: Choose Next Simulation block

- Check Regime Comparison Table → choose untested combination
- **Do not replicate** previous block unless motivated (testing knowledge transfer)

### STEP 3: Update Working Memory

Update `{config}_memory.md`:

- Update Knowledge Base with confirmed principles
- Add row to Regime Comparison Table
- Replace Previous Block Summary with **short summary** (2-3 lines, NOT individual iterations)
- Clear "Iterations This Block" section
- Write hypothesis for next block

---

## Working Memory Structure

```markdown
# Working Memory

## Knowledge Base (accumulated across all blocks)

### Regime Comparison Table

| Block | Regime                  | E/I | n_frames | n_neurons | n_types | eff_rank | Best R² | Optimal lr_W | Optimal L1 | Key finding |
| ----- | ----------------------- | --- | -------- | --------- | ------- | -------- | ------- | ------------ | ---------- | ----------- |
| 1     | chaotic, Dale_law=False | -   | 10000    | 100       | 1       | 31-35    | 1.000   | 8E-3         | 1E-5       | ...         |

### Established Principles

[Confirmed patterns that apply across regimes]

### Open Questions

[Patterns needing more testing, contradictions]

---

## Previous Block Summary (Block N-1)

[Short summary only - NOT individual iterations. Example:
"Block 1 (chaotic, Dale_law=False): Best R²=1.000 at lr_W=8E-3.
Key finding: lr_W optimal range 4E-3 to 8E-3."]

---

## Current Block (Block N)

### Block Info

Simulation: connectivity_type=X, Dale_law=Y, ...
Iterations: M to M+n_iter_block

### Hypothesis

[Prediction for this block, stated before running]

### Iterations This Block

[Current block iterations only — cleared at block boundary]

### Emerging Observations

[Running notes on what's working/failing]
**CRITICAL: This section must ALWAYS be at the END of memory file. When adding new iterations, insert them BEFORE this section.**
```

---

## Knowledge Base Guidelines

### What to Add to Established Principles

Examples:

- ✓ "Constrained connectivity needs lower lr_W" (causal, generalizable)
- ✓ "L1 > 1e-04 fails for low_rank" (boundary condition)
- ✓ "Effective rank < 15 requires factorization=True" (theoretical link)
- ✗ "lr_W=0.01 worked in Block 4" (too specific)
- ✗ "Block 3 converged" (not a principle)

### Scientific Method (CRITICAL)

**Repeatability:**

- Each iteration is a single training run with stochastic initialization
- Results may vary between runs with identical config
- A "failed" run may succeed on retry; a "converged" run may fail

**Evidence hierarchy:**

| Level            | Criterion                              | Action                 |
| ---------------- | -------------------------------------- | ---------------------- |
| **Established**  | Consistent across 3+ iterations/blocks | Add to Principles      |
| **Tentative**    | Observed 1-2 times                     | Add to Open Questions  |
| **Contradicted** | Conflicting evidence                   | Note in Open Questions |

### What to Add to Open Questions

Examples:

- Patterns needing more testing
- Contradictions between blocks
- Theoretical predictions not yet verified

---

## Theoretical Background

### GNN Architecture (Signal_Propagation)

The model learns neural dynamics du/dt using a graph neural network:

```
du/dt = lin_phi(u, a) + W @ lin_edge(u, a)
```

**Components:**

- `lin_edge` (MLP): message function on edges, transforms source neuron activity
- `lin_phi` (MLP): node update function, computes local dynamics
- `W`: learnable connectivity matrix (n_neurons × n_neurons)
- `a`: learnable node embeddings (n_neurons × embedding_dim)

**Forward pass:**

1. For each edge (j→i): message = W[i,j] × lin_edge(u_j, a_j)
2. Aggregate messages: msg_i = Σ_j W[i,j] × lin_edge(u_j, a_j)
3. Update: du/dt_i = lin_phi(u_i, a_i) + msg_i

**Low-rank factorization:**

When `low_rank_factorization=True`: W = W_L @ W_R where W_L ∈ ℝ^(n×r), W_R ∈ ℝ^(r×n)

**Node embeddings for heterogeneity:**

The embedding vector `a_i` allows each neuron to have different dynamics parameters:

- When `n_neuron_types > 1`: embeddings are learnable (requires `lr_emb`)
- When `n_neuron_types = 1`: embeddings are fixed (all neurons identical)
- Embeddings are concatenated with activity: lin_phi(u_i, a_i) and lin_edge(u_j, a_j)
- This allows the MLPs to learn neuron-type-specific transfer functions

### Training Loss and Regularization

**Prediction loss:**

```
L_pred = ||du/dt_pred - du/dt_true||₂
```

**Key regularization terms:**

- `coeff_W_L1`: L1 on W (sparsity). Range: 1E-6 to 1E-4
- `coeff_edge_diff`: enforces monotonicity of lin_edge output (positive: higher u → higher output)
  - Computed as: relu(msg(u) - msg(u+δ))·coeff for sampled u values
  - Stabilizes message function, prevents oscillating gradients

**Learning rates:**

- `learning_rate_W_start` (lr_W): learning rate for connectivity W
- `learning_rate_start` (lr): learning rate for lin_edge and lin_phi
- `learning_rate_embedding_start` (lr_emb): learning rate for node embeddings a

**Total loss:**

```
L = L_pred + coeff_W_L1·||W||₁ + coeff_edge_diff·L_edge_diff + ...
```

### Spectral Radius

- ρ(W) < 1: activity decays → harder to constrain W
- ρ(W) ≈ 1: edge of chaos → rich dynamics → good recovery
- ρ(W) > 1: unstable

### Effective Rank

- High (30+): full W recoverable
- Low (<15): only subspace identifiable → need factorization
- **Block 6 finding**: effective_rank is the primary predictor of achievable R²
  - effective_rank=20 → R²≈0.998 achievable
  - effective_rank=10 → R²≈0.92 ceiling (regardless of training params)

### Low-rank Connectivity

- W = W_L @ W_R constrains solution space
- Without factorization: spurious full-rank solutions

### Learning Rates

- lr_W:lr ratio matters (typically 20:1 to 50:1)
- Too fast φ learning → noisy W gradients

\end{verbatim}
\end{small}

\section{Memory File}
\label{app:memory}

\begin{small}
\begin{verbatim}
# Working Memory: signal_chaotic_1_Claude

## Knowledge Base (accumulated across all blocks)

### Regime Comparison Table
| Block | Regime | eff_rank | Best R² | Optimal lr_W | Optimal L1 | Key finding |
|-------|--------|----------|---------|--------------|------------|-------------|
| 1 | chaotic, Dale_law=False, n_frames=10000 | 31-35 | 1.000 | 4E-3 to 4E-2 | ≤5E-4 | extremely robust; low_rank_factorization=True catastrophic |
| 2 | chaotic, Dale_law=True, n_frames=10000 | 10 | 0.913 | 5E-2 to 1.2E-1 | ≤1E-5 | R² ceiling ~0.91 due to low eff_rank; needs 10-30x higher lr_W |
| 3 | low_rank (r=20), Dale_law=False | 6 (stochastic) | 0.886* | 8E-2 to 1E-1 | ≤1E-5 | *not reproducible; eff_rank=6 in 15/16 runs creates R²<0.4 ceiling |
| 4 | low_rank (r=50), Dale_law=False + factorization | 7 | 0.188 | 1.5E-1 | 1E-6 | factorization=True helps (10x), but eff_rank=7 hard ceiling |
| 5 | chaotic, Dale_law=False, n_frames=5000 | 20 | 0.924 | 2E-2 to 3E-2 | 1E-5 | n_frames reduction lowers eff_rank; needs 5-7x higher lr_W |
| 6 | chaotic, Dale_law=False, n_frames=2500 | 6* | 0.568 | N/A | 1E-5 | *94% eff_rank=6; below minimum data threshold; 0/16 converged |
| 7 | chaotic, Dale_law=False, n_frames=7500 | 25-29 | 0.989 | 3E-2 to 1.3E-1 | ≤5E-5 | intermediate data: R²=0.989 achievable; lr_W 40x range |
| 8 | chaotic, Dale_law=True, n_frames=10000 | 10 (94%), 27 (6%) | 0.995 | 8E-2 to 1.5E-1 | 1E-6 | eff_rank highly stochastic; eff_rank=27 → R²=0.995; eff_rank=10 → R²≤0.945 |
| 9 | low_rank (r=20), Dale_law=True, n_frames=10000 | 16 | 0.540 | TBD | 1E-6 | single iteration; eff_rank=16 (better than low_rank+Dale_law=False) |

### Established Principles
1. chaotic full-rank connectivity: lr_W has 10x robust range (4E-3 to 4E-2)
2. L1 regularization boundary: regime-dependent (Dale_law=False: ≤5E-4, Dale_law=True: ≤1E-5, low_rank: ≤1E-6)
3. low_rank_factorization=True: fails on chaotic connectivity; HELPS on low_rank connectivity (when factorization rank ≥ connectivity rank)
4. batch_size (8, 16, 32) does not affect convergence in any tested regime
5. Dale_law=True: eff_rank stochastic (Block 2: eff_rank≈10, Block 8: eff_rank=10-27) - NOT deterministic
6. effective_rank determines R² ceiling: eff_rank=6 → R²≈0.15-0.57; eff_rank=10 → R²≈0.91-0.95; eff_rank=20 → R²≈0.92; eff_rank≥25 → R²≈0.99; eff_rank=32 → R²=1.0
7. Dale_law=True requires 10-30x higher lr_W than Dale_law=False
8. lr/lr_W ratio matters: optimal ~1000:1 to 2000:1 for Dale_law=True; 10000:1 for low_rank
9. low_rank connectivity produces lower eff_rank than chaotic; high stochasticity in eff_rank
10. lr_W optimal range is regime-dependent: chaotic ~4E-3, Dale_law=True ~8E-2 to 1.5E-1, low_rank ~1.5E-1
11. connectivity_rank does NOT meaningfully affect eff_rank (rank=50 gave eff_rank=7, similar to rank=20 with eff_rank=6)
12. n_frames affects eff_rank: n_frames=2500 → eff_rank≈6; n_frames=5000 → eff_rank≈20; n_frames=7500 → eff_rank≈27; n_frames=10000 → eff_rank≈32
13. lower eff_rank requires higher lr_W: eff_rank=20 (n_frames=5000) needs 5-7x higher lr_W (2E-2 to 3E-2 vs 4E-3)
14. **minimum n_frames for chaotic regime ≈ 5000** (n_frames=2500 fails 94% with eff_rank=6)
15. n_frames=7500 provides good balance: eff_rank≈27, R²=0.989 achievable, 40x lr_W range
16. **eff_rank=10 ceiling can reach R²=0.945** with optimized params (lr_W=1.5E-1, lr=1E-4, L1=1E-6)

### Open Questions
- why does Dale_law=True sometimes give eff_rank≈10 (94%) vs eff_rank=27 (6%)? stochasticity in initialization?
- does n_neuron_types>1 change parameter sensitivity? (test after iter 1024)
- can low_rank + Dale_law=True achieve convergence? (initial result: eff_rank=16, R²=0.540 - promising but inconclusive)
- does eff_rank=16 regime have higher R² ceiling than eff_rank=10?

---

## Previous Block Summary

**Block 9** (low_rank r=20, Dale_law=True, n_frames=10000): 0/1 converged, Best R²=0.540 (iter 128).
Key finding: eff_rank=16 (better than Block 3's eff_rank=6 for low_rank+Dale_law=False); single iteration - need more data.

---

## Current Block (Block 10)

### Block Info
Simulation: connectivity_type=low_rank, connectivity_rank=20, Dale_law=True, n_frames=10000, n_neurons=100, n_types=1
Iterations: 129 to 144

### Hypothesis
Continuing low_rank + Dale_law=True exploration:
1. Block 9 showed eff_rank=16 (better than expected), R²=0.540 with lr_W=1.5E-1, L1=1E-6
2. eff_rank=16 should have higher R² ceiling than eff_rank=10 (~0.945) - estimate ~0.85-0.95
3. May need to tune lr_W (try 8E-2 to 2E-1 range) and lr (try 5E-5 to 2E-4)
4. low_rank_factorization=True is set - this helps for low_rank connectivity
5. Prediction: with optimization, R²≈0.7-0.9 achievable

### Iterations This Block

## Iter 129: partial
Node: id=129, parent=root
Config: lr_W=8E-2, lr=1E-4, L1=1E-6, low_rank_factorization=T
Metrics: connectivity_R2=0.612
Observation: eff_rank=16 consistent; baseline for block

## Iter 130: partial
Node: id=130, parent=129
Config: lr_W=1.5E-1, lr=1E-4, L1=1E-6, low_rank_factorization=T
Metrics: connectivity_R2=0.539
Mutation: lr_W: 8E-2 -> 1.5E-1
Observation: R² dropped with higher lr_W; suggests lower lr_W better for eff_rank=16

## Iter 131: partial
Node: id=131, parent=130
Config: lr_W=1E-1, lr=1E-4, L1=1E-6, low_rank_factorization=T
Metrics: connectivity_R2=0.584
Mutation: lr_W: 1.5E-1 -> 1E-1 (0.67x)
Observation: R²=0.584 intermediate; confirms lower lr_W trend

## Iter 132: partial
Node: id=132, parent=131
Config: lr_W=6E-2, lr=1E-4, L1=1E-6, low_rank_factorization=T
Metrics: connectivity_R2=0.633
Mutation: lr_W: 1E-1 -> 6E-2 (0.6x)
Observation: new best R²=0.633; trend: 1.5E-1→0.539, 1E-1→0.584, 8E-2→0.612, 6E-2→0.633

## Iter 133: partial
Node: id=133, parent=132
Config: lr_W=4E-2, lr=1E-4, L1=1E-6, low_rank_factorization=T
Metrics: connectivity_R2=0.657, test_R2=0.301, test_pearson=0.299
Mutation: lr_W: 6E-2 -> 4E-2 (0.67x)
Observation: new best R²=0.657; trend continues: lower lr_W = higher R²

## Iter 134: partial
Node: id=134, parent=133
Config: lr_W=3E-2, lr=1E-4, L1=1E-6, low_rank_factorization=T
Metrics: connectivity_R2=0.668, test_R2=0.314, test_pearson=0.335
Mutation: lr_W: 4E-2 -> 3E-2 (0.75x)
Observation: new best R²=0.668; trend continues but diminishing returns (~0.011 improvement vs ~0.024 prior)

## Iter 135: partial
Node: id=135, parent=134
Config: lr_W=2E-2, lr=1E-4, L1=1E-6, low_rank_factorization=T
Metrics: connectivity_R2=0.669, test_R2=0.300, test_pearson=0.309
Mutation: lr_W: 3E-2 -> 2E-2 (0.67x)
Observation: R²=0.669 flat; lr_W optimum found at 2E-2 to 3E-2

## Iter 136: partial
Node: id=136, parent=135
Config: lr_W=2E-2, lr=5E-5, L1=1E-6, low_rank_factorization=T
Metrics: connectivity_R2=0.689, test_R2=0.294, test_pearson=0.128
Mutation: lr: 1E-4 -> 5E-5 (0.5x)
Observation: new best R²=0.689! lr=5E-5 helps; lr/lr_W ratio 400:1

## Iter 137: partial
Node: id=137, parent=136
Config: lr_W=2E-2, lr=2E-5, L1=1E-6, low_rank_factorization=T
Metrics: connectivity_R2=0.681, test_R2=0.322, test_pearson=0.377
Mutation: lr: 5E-5 -> 2E-5 (0.4x)
Observation: R²=0.681 vs iter 136's 0.689; lr=2E-5 slightly worse; lr optimum ~5E-5
Next: parent=137 (UCB=2.802); try lr=8E-5 to explore higher lr

### Emerging Observations

- eff_rank=16 confirmed for low_rank+Dale_law regime (10 consecutive runs: 128-137)
- **lr_W optimum found**: 2E-2 to 3E-2 (trend plateaued at R²≈0.669)
- Full lr_W trend: 1.5E-1→0.539, 1E-1→0.584, 8E-2→0.612, 6E-2→0.633, 4E-2→0.657, 3E-2→0.668, 2E-2→0.669
- eff_rank=16 prefers much lower lr_W than eff_rank=10 (2E-2 vs 8E-2 to 1.5E-1)
- **lr optimum found**: lr=5E-5 (R²=0.689) > lr=2E-5 (R²=0.681) > lr=1E-4 (R²=0.669)
- optimal lr/lr_W ratio for eff_rank=16: ~400:1 (lr=5E-5 / lr_W=2E-2)
- next: try lr=8E-5 to confirm lr=5E-5 is optimal (explore higher direction)


\end{verbatim}
\end{small}


% \paragraph{Epistemic metrics.}
% To characterize system-level epistemic behavior beyond task performance, we introduce three minimal metrics computed from existing logs and memory files. \emph{Hypothesis persistence} measures the duration over which explicit hypotheses or principles remain active in structured memory across iterations or blocks, quantifying cumulative induction versus unstable exploration. \emph{Diagnostic experiment fraction} measures the proportion of runs explicitly devoted to hypothesis testing, boundary probing, or failure disambiguation, as identified by logged exploration modes, capturing abductive and falsificatory behavior rather than pure exploitation. Finally, \emph{regime transfer efficiency} measures the reduction in iterations required to reach a predefined convergence criterion when entering new simulation regimes, quantifying the reuse of previously established constraints and accumulated knowledge. Together, these metrics assess how effectively the system acquires, tests, revises, and transfers knowledge over extended experimental sequences.

\appendix
\section{Epistemic Metrics}
\label{appendix:epistemic_metrics}

We introduce \emph{epistemic metrics} to quantify how the system acquires, tests, revises, and transfers knowledge across experimental blocks. These metrics characterize system-level behaviors such as hypothesis refinement, diagnostic experimentation, regime differentiation, and the efficient transfer of previously established principles when entering new regions of the search space. Importantly, these metrics do not measure task performance alone, but the structure and dynamics of the exploration process itself.

\subsection{Reasoning Mode Definitions}

We classify reasoning events into twelve epistemic modes, each capturing a distinct cognitive operation:

\begin{description}
    \item[Induction] Pattern extraction from accumulated observations. Measured by the number of observations required before a generalizable pattern is articulated.
    \item[Abduction] Generation of causal hypotheses to explain unexpected observations. Typically triggered by performance anomalies or regime changes.
    \item[Deduction] Derivation of testable predictions from hypotheses. Enables hypothesis validation through targeted experiments.
    \item[Falsification] Rejection of hypotheses based on contradictory evidence. Critical for pruning the hypothesis space and refining understanding.
    \item[Boundary] Systematic probing of parameter limits to establish operating ranges. Identifies safe vs.\ failure regions.
    \item[Analogy/Transfer] Application of knowledge from one regime to another. Success rate indicates generalizability of learned principles.
    \item[Meta-reasoning] Reflection on and adaptation of exploration strategy. Triggered when current approaches prove ineffective.
    \item[Regime] Identification of qualitatively distinct operating phases. Enables context-appropriate hypothesis generation.
    \item[Uncertainty] Recognition of stochastic or irreducible variability. Informs confidence bounds and exploration intensity.
    \item[Causal] Construction of multi-step mechanistic chains linking variables. Represents deep structural understanding.
    \item[Predictive] Formulation of quantitative rules enabling \emph{a priori} predictions. Highest form of operational knowledge.
    \item[Constraint] Derivation of parameter relationships that must hold for success. Enables principled configuration.
\end{description}

\subsection{Primary Metrics}

\subsubsection{Hypothesis Turnover Rate (HTR)}

The rate at which hypotheses are generated and subsequently falsified or refined:
\begin{equation}
    \text{HTR} = \frac{N_{\text{falsified}} + N_{\text{refined}}}{N_{\text{abduction}}}
\end{equation}
where $N_{\text{falsified}}$ is the count of falsification events, $N_{\text{refined}}$ the count of hypothesis refinements, and $N_{\text{abduction}}$ the count of abductive hypotheses generated. High HTR indicates active hypothesis testing; HTR $\approx 1$ suggests efficient single-shot hypothesis validation; HTR $> 1$ indicates iterative refinement cycles.

\subsubsection{Deductive Accuracy (DA)}

The fraction of deductive predictions confirmed by subsequent experiments:
\begin{equation}
    \text{DA} = \frac{N_{\text{deduction}}^{\checkmark}}{N_{\text{deduction}}^{\checkmark} + N_{\text{deduction}}^{\times}}
\end{equation}
where $N_{\text{deduction}}^{\checkmark}$ counts confirmed predictions and $N_{\text{deduction}}^{\times}$ counts falsified predictions. DA reflects the quality of the underlying hypotheses and the system's predictive calibration.

\subsubsection{Transfer Success Rate (TSR)}

The effectiveness of cross-regime knowledge transfer:
\begin{equation}
    \text{TSR} = \frac{N_{\text{transfer}}^{\text{success}} + 0.5 \cdot N_{\text{transfer}}^{\text{partial}}}{N_{\text{transfer}}^{\text{total}}}
\end{equation}
where transfer attempts are classified as success ($\checkmark$), partial ($\sim$), or failure ($\times$). TSR quantifies the generalizability of learned principles across different experimental regimes.

\subsubsection{Causal Depth (CD)}

The average length of mechanistic causal chains constructed:
\begin{equation}
    \text{CD} = \frac{1}{N_{\text{causal}}} \sum_{i=1}^{N_{\text{causal}}} L_i
\end{equation}
where $L_i$ is the number of variables in causal chain $i$. Higher CD indicates deeper mechanistic understanding beyond surface correlations.

\subsubsection{Epistemic Density (ED)}

The concentration of reasoning events per iteration:
\begin{equation}
    \text{ED} = \frac{N_{\text{events}}}{N_{\text{iterations}}}
\end{equation}
where $N_{\text{events}}$ is the total count of classified reasoning events. ED captures the intensity of epistemic activity; higher values indicate richer reasoning per experimental step.

\subsection{Structural Metrics}

\subsubsection{Falsification Latency (FL)}

The average number of iterations between hypothesis generation and falsification:
\begin{equation}
    \text{FL} = \frac{1}{N_{\text{falsified}}} \sum_{j=1}^{N_{\text{falsified}}} (t_j^{\text{falsify}} - t_j^{\text{hypothesis}})
\end{equation}
Short FL indicates rapid hypothesis testing; long FL may indicate either robust hypotheses or insufficient testing.

\subsubsection{Induction Consolidation Time (ICT)}

The number of observations required before pattern articulation:
\begin{equation}
    \text{ICT} = \frac{1}{N_{\text{induction}}} \sum_{k=1}^{N_{\text{induction}}} n_k^{\text{obs}}
\end{equation}
where $n_k^{\text{obs}}$ is the observation count preceding induction event $k$. Lower ICT indicates faster pattern recognition.

\subsubsection{Regime Sensitivity Index (RSI)}

The fraction of block transitions that trigger regime recognition:
\begin{equation}
    \text{RSI} = \frac{N_{\text{regime}}}{N_{\text{blocks}} - 1}
\end{equation}
RSI $= 1$ indicates perfect regime awareness; lower values suggest failure to recognize qualitative changes in the problem structure.

\subsubsection{Meta-Reasoning Trigger Rate (MTR)}

The frequency of strategy-level adaptation:
\begin{equation}
    \text{MTR} = \frac{N_{\text{meta}}}{N_{\text{blocks}}}
\end{equation}
Higher MTR indicates more frequent strategic reflection, typically triggered by persistent failures or unexpected successes.

\subsection{Graph-Theoretic Metrics}

The epistemic timeline can be represented as a directed graph $G = (V, E)$ where vertices $V$ are reasoning events and edges $E$ represent causal or temporal dependencies.

\subsubsection{Reasoning Chain Length (RCL)}

The average path length in the epistemic graph:
\begin{equation}
    \text{RCL} = \frac{1}{|P|} \sum_{p \in P} |p|
\end{equation}
where $P$ is the set of maximal directed paths and $|p|$ is path length. Longer chains indicate sustained reasoning threads.

\subsubsection{Cross-Block Connectivity (CBC)}

The ratio of cross-block to within-block edges:
\begin{equation}
    \text{CBC} = \frac{|E_{\text{cross}}|}{|E_{\text{within}}|}
\end{equation}
Higher CBC indicates stronger integration of knowledge across experimental phases.

\subsubsection{Falsification Fan-In (FFI)}

The average number of hypotheses rejected by each falsification event:
\begin{equation}
    \text{FFI} = \frac{|\{e \in E : \text{target}(e) \in V_{\text{falsification}}\}|}{|V_{\text{falsification}}|}
\end{equation}
FFI $> 1$ indicates efficient falsification that eliminates multiple competing hypotheses simultaneously.

\subsection{Temporal Dynamics}

\subsubsection{Epistemic Phase Transition}

We identify phase transitions in reasoning complexity by tracking the emergence of advanced modes (Causal, Predictive, Constraint) over time. Define the \emph{epistemic maturity} at iteration $t$:
\begin{equation}
    M(t) = \frac{N_{\text{causal}}(t) + N_{\text{predictive}}(t) + N_{\text{constraint}}(t)}{N_{\text{events}}(t)}
\end{equation}
The first iteration where $M(t) > 0$ marks the transition from exploratory to explanatory reasoning.

\subsubsection{Knowledge Crystallization Rate (KCR)}

The rate at which predictive rules emerge from accumulated experience:
\begin{equation}
    \text{KCR} = \frac{N_{\text{predictive}}}{N_{\text{iterations}}} \times 100
\end{equation}
expressed as predictive rules per 100 iterations. Higher KCR indicates efficient distillation of experience into actionable knowledge.

\subsection{Computed Metrics for signal\_chaotic\_1\_Claude}

Table~\ref{tab:epistemic_metrics} presents the epistemic metrics computed from 149 iterations across 11 experimental blocks.

\begin{table}[htbp]
\centering
\caption{Epistemic metrics for the signal\_chaotic\_1\_Claude experiment.}
\label{tab:epistemic_metrics}
\begin{tabular}{lrp{6cm}}
\toprule
\textbf{Metric} & \textbf{Value} & \textbf{Interpretation} \\
\midrule
\multicolumn{3}{l}{\emph{Primary Metrics}} \\
Hypothesis Turnover Rate (HTR) & 2.43 & Active iterative refinement \\
Deductive Accuracy (DA) & 0.71 & 71\% of predictions confirmed \\
Transfer Success Rate (TSR) & 0.35 & Limited cross-regime generalization \\
Causal Depth (CD) & 3.6 & Multi-step mechanistic chains \\
Epistemic Density (ED) & 0.72 & 0.72 reasoning events per iteration \\
\midrule
\multicolumn{3}{l}{\emph{Structural Metrics}} \\
Falsification Latency (FL) & 4.2 iter & Rapid hypothesis testing \\
Induction Consolidation Time (ICT) & 6.8 obs & Moderate pattern recognition speed \\
Regime Sensitivity Index (RSI) & 0.70 & 7/10 block transitions recognized \\
Meta-Reasoning Trigger Rate (MTR) & 0.55 & Strategy adaptation in 6/11 blocks \\
\midrule
\multicolumn{3}{l}{\emph{Graph-Theoretic Metrics}} \\
Total Edges & 53 & Causal relationships identified \\
Within-Block Edges & 44 & Local reasoning chains \\
Cross-Block Edges & 9 & Knowledge transfer links \\
Cross-Block Connectivity (CBC) & 0.20 & Moderate integration \\
\midrule
\multicolumn{3}{l}{\emph{Temporal Dynamics}} \\
First Causal Chain & Iter 48 & Epistemic phase transition \\
Knowledge Crystallization Rate (KCR) & 2.01 & 3 predictive rules / 149 iterations \\
\bottomrule
\end{tabular}
\end{table}

\subsection{Mode Distribution}

The distribution of reasoning events by mode reveals the epistemic profile of the exploration:

\begin{table}[htbp]
\centering
\caption{Distribution of reasoning events by epistemic mode.}
\label{tab:mode_distribution}
\begin{tabular}{lrrp{5cm}}
\toprule
\textbf{Mode} & \textbf{Count} & \textbf{\%} & \textbf{Role} \\
\midrule
Falsification & 20 & 18.5\% & Hypothesis elimination \\
Deduction & 17 & 15.7\% & Prediction generation \\
Boundary & 12 & 11.1\% & Parameter limit probing \\
Induction & 13 & 12.0\% & Pattern recognition \\
Analogy/Transfer & 10 & 9.3\% & Cross-regime application \\
Abduction & 10 & 9.3\% & Hypothesis generation \\
Regime & 7 & 6.5\% & Phase identification \\
Meta-reasoning & 6 & 5.6\% & Strategy adaptation \\
Uncertainty & 5 & 4.6\% & Stochasticity awareness \\
Causal & 5 & 4.6\% & Mechanistic modeling \\
Predictive & 3 & 2.8\% & Quantitative rules \\
Constraint & 1 & 0.9\% & Parameter relationships \\
\midrule
\textbf{Total} & \textbf{108} & \textbf{100\%} & \\
\bottomrule
\end{tabular}
\end{table}

\subsection{Interpretation}

The epistemic metrics reveal several key characteristics of the LLM-driven exploration:

\paragraph{Hypothesis-Driven Exploration.} The high HTR (2.43) indicates that hypotheses undergo multiple refinement cycles rather than single-shot validation. Combined with DA = 0.71, this suggests calibrated but not overconfident prediction.

\paragraph{Limited Transfer.} TSR = 0.35 reflects the challenge of cross-regime generalization. Most knowledge learned in one block required adaptation when applied to structurally different regimes (e.g., Dale's law, low-rank constraints).

\paragraph{Emergent Causal Understanding.} The appearance of causal chains at iteration 48 marks a phase transition from parameter search to mechanistic modeling. The average causal depth of 3.6 indicates multi-step reasoning beyond pairwise correlations.

\paragraph{Balanced Reasoning Profile.} The mode distribution shows a roughly balanced profile between constructive (Induction, Abduction, Deduction) and critical (Falsification, Boundary) reasoning, with meta-cognitive modes (Meta-reasoning, Regime, Uncertainty) providing adaptive control.

\paragraph{Knowledge Crystallization.} Despite 149 iterations, only 3 predictive rules emerged (KCR = 2.01 per 100 iterations), suggesting that converting experimental experience into generalizable quantitative principles remains challenging.


\section{Cumulative scientific inference}

\renewcommand{\arraystretch}{1.3}

\subsection*{Table A1: Induction (observations $\rightarrow$ pattern)}

\begin{table}[h!]
\centering
\small
\begin{tabularx}{\textwidth}{XX}
\toprule
\textbf{Single-shot} & \textbf{Cumulative} \\
\midrule
``lr\_W=8E-3 works for this config'' (iter 3)
&
``lr\_W:lr ratio scales with constraint complexity: 40:1 chaotic, 320:1 low\_rank, 800:1 Dale'' (experiments 1--7)
\\
\midrule
``effective\_rank=32, R$^2$=1.0'' (iter 2)
&
``effective\_rank $\geq$ 20 $\Rightarrow$ R$^2 \approx$ 0.999; effective\_rank $<$ 15 $\Rightarrow$ R$^2 \leq$ 0.92'' (experiments 1--8)
\\
\bottomrule
\end{tabularx}
\end{table}

\subsection*{Table A2: Abduction (observation $\rightarrow$ explanatory hypothesis)}

\begin{table}[h!]
\centering
\small
\begin{tabularx}{\textwidth}{XX}
\toprule
\textbf{Single-shot} & \textbf{Cumulative} \\
\midrule
``R$^2$ jumped 0.16 to 0.89'' (iter 125)
&
``effective\_rank jumped 4$\rightarrow$7; this causes R$^2$ improvement---12 prior failures at eff\_rank=4 prove it'' (iters 113--127)
\\
\midrule
``n\_frames=10000 insufficient for Dale'' (experiment 3)
&
``n\_frames controls effective\_rank ceiling; 20k enables eff\_rank 28--36 vs 10'' (experiments 3, 7)
\\
\bottomrule
\end{tabularx}
\end{table}

\subsection*{Table A3: Deduction (hypothesis $\rightarrow$ prediction)}

\begin{table}[h!]
\centering
\small
\begin{tabularx}{\textwidth}{XX}
\toprule
\textbf{Single-shot} & \textbf{Cumulative} \\
\midrule
``lr\_W=8E-3 worked $\Rightarrow$ try 10E-3'' (iter 4)
&
``ratio 80:1 optimal experiment 1; experiment 2 needs 4$\times$ $\Rightarrow$ predict lr\_W=16E-3, lr=5E-5'' (iter 29, verified)
\\
\midrule
``maybe n\_frames=30000 helps'' (iter 124)
&
``$\mathcal{H}_1$: n\_frames$\uparrow$ $\Rightarrow$ eff\_rank$\uparrow$; $\mathcal{H}_2$: eff\_rank$\uparrow$ $\Rightarrow$ R$^2\uparrow$; $\therefore$ predict n\_frames=30k helps'' (iter 124, falsified)
\\
\bottomrule
\end{tabularx}
\end{table}

\subsection*{Table A4: Falsification (prediction vs.\ outcome $\rightarrow$ reject/refine)}

\begin{table}[h!]
\centering
\small
\begin{tabularx}{\textwidth}{XX}
\toprule
\textbf{Single-shot} & \textbf{Cumulative} \\
\midrule
``low\_rank unlearnable, give up'' (iters 17--25)
&
``falsified: low\_rank=20 unlearnable; refined: learnable with ratio$\geq$320:1'' (breakthrough iter 29)
\\
\midrule
``n\_frames=30000 didn't help'' (iter 124)
&
``falsified: n\_frames always boosts eff\_rank; refined: only when connectivity\_rank $>$ threshold'' (eff\_rank=4 unchanged)
\\
\bottomrule
\end{tabularx}
\end{table}
\end{document}


